\documentclass[../DoAn.tex]{subfiles}
\begin{document}

Phần tài liệu tham khảo của đồ án ưu tiên các bài báo khảo sát, bài báo tạp chí, sách nền tảng và tài liệu kỹ thuật liên quan trực tiếp tới Coverage Path Planning, lập kế hoạch chuyển động, môi trường động, ràng buộc năng lượng, thuật toán tìm đường và mô hình truyền động của robot di động. Các nguồn dạng slide bài giảng, trang tổng hợp không chính thống hoặc tài liệu không có thông tin xuất bản rõ ràng không được sử dụng làm tài liệu chính.

Trong số các tài liệu được dùng làm nền tảng, bài khảo sát về Coverage Path Planning cho robot di động trong môi trường động là tài liệu quan trọng nhất vì cung cấp bức tranh tổng quan về bài toán CPP, các nhóm thuật toán, môi trường tĩnh và động, vật cản động, chỉ số đánh giá và các hướng nghiên cứu gần đây. Các tài liệu khảo sát kinh điển về CPP được dùng để đặt bài toán vào bối cảnh chung của robot tự hành. Các tài liệu về lập kế hoạch chuyển động, thuật toán Dijkstra và robot di động được dùng để bổ sung nền tảng cho phần tìm đường trên đồ thị có trọng số, trạng thái có hướng và thảo luận về cấu trúc truyền động.

Các tài liệu kỹ thuật từ các nền tảng robot thương mại được sử dụng với vai trò minh họa cho sự khác biệt giữa nhóm robot có hướng thân rõ ràng và nhóm robot toàn hướng. Đây là nhóm tài liệu bổ trợ cho lập luận mô hình hóa, không phải phần thay thế cho tài liệu nghiên cứu thuật toán. Cách sử dụng này phù hợp với phạm vi của đồ án: báo cáo tập trung vào mô hình thuật toán và mô phỏng trong ngành Khoa học máy tính, còn các chi tiết phần cứng, cảm biến và triển khai thực địa chỉ được dùng ở mức bối cảnh kỹ thuật.

\nocite{jayalakshmi2025comprehensive}
\nocite{choset2001coverage}
\nocite{galceran2013survey}
\nocite{shnaps2016online}
\nocite{tan2021comprehensive}
\nocite{jensen2021near}
\nocite{ramesh2024anytime}
\nocite{cai2020mobile}
\nocite{latombe1991robot}
\nocite{lavalle2006planning}
\nocite{dijkstra1959note}
\nocite{siegwart2011autonomous}
\nocite{palroboticsTiagoBase2026}
\nocite{palroboticsOmniBase2026}
\nocite{clearpathHuskyA3002026}
\nocite{revroboticsMaxSwerve2026}

\end{document}
